\subsection{Audit of the assertion that changing the complex structure merges any two packets}

The notebook, P000702, third line from the bottom before the final diagram, says
\ru{Любые два неприв. предст. эквивал. относит. нек-рой рим. пов. (вх. в какой-то пакет).}
The scan confirms the two universal quantifiers; the word abbreviated as
\ru{неприв.} is overwritten but is the diplomatic transcription's reading.
The definition earlier on the same page makes equivalence mean
\(E_{\rho_1}\cong E_{\rho_2}\) as holomorphic bundles.  The assertion therefore
includes the following rank-one case, in which it is false.  The author's own
survey explicitly gives uniqueness of the unitary representation in a packet
as an exercise \citep[Ch.~II, \S2, Exercise~10, p.~204]{tyurin1978periods}.
This is an audit
of a literal assertion, not the development of a C or D problem.

\begin{proposition}[Distinct unitary characters remain in distinct packets]\ours{}
Let \(\Sigma\) be a connected, closed, oriented surface of genus \(g\geq1\),
and let \(\chi_1,\chi_2:\pi_1(\Sigma)\to U(1)\) be two distinct characters.
For every complex structure \(J\) compatible with the orientation, the
holomorphic line bundles obtained from the two unitary flat bundles on
\(X=(\Sigma,J)\) are nonisomorphic.  In particular, the notebook's assertion
about arbitrary pairs of irreducible representations is false.
\end{proposition}

\begin{proof}
Write \(L_i\) for the flat line bundle with monodromy \(\chi_i\), with its
parallel Hermitian metric and holomorphic structure given by the
\((0,1)\)-part of the flat connection.  Suppose that
\(\Phi:L_1\to L_2\) is a holomorphic bundle isomorphism.  It is a nowhere-zero
holomorphic section \(s\) of the unitary flat line bundle
\(L=L_2\otimes L_1^{-1}\), whose monodromy is
\(\delta=\chi_2\chi_1^{-1}\).

On any simply connected coordinate neighbourhood choose a parallel frame
\(e\) of norm one for \(L\).  Because the \((0,1)\)-connection defines the
holomorphic structure, \(e\) is holomorphic.  Thus \(s=f e\) with \(f\)
holomorphic and nowhere zero.  The real function
\(u=\log\lVert s\rVert=\log|f|\) is harmonic on that neighbourhood.  Unitary
transition functions make \(u\) a globally defined harmonic function on the
compact connected Riemann surface \(X\); hence \(u\) is constant.  Each
holomorphic function \(f\) therefore has constant modulus and is constant.
Consequently \(s\) is a nonzero parallel section of \(L\).  Parallel
transport around every loop fixes its nonzero value, so \(\delta=1\) and
\(\chi_1=\chi_2\), a contradiction.

For an explicit pair, use the presentation in the notebook.  Take
\(\chi_1=1\), let \(\chi_2(a_1)=e^{i\theta}\) for
\(\theta\notin2\pi\mathbb Z\), and send every other standard generator to
\(1\).  The surface relation holds because \(U(1)\) is abelian.  Both are
one-dimensional complex representations and therefore irreducible.  The
argument applies to every compatible \(J\), including genus \(g\geq2\)
if that was the intended range.  Even allowing a change of marking cannot
turn the nontrivial character into the trivial character.
\end{proof}

\begin{remark}[Independent check]
The proof was checked again by working on the universal cover.  A
holomorphic isomorphism there is multiplication by a nowhere-zero
holomorphic function \(F\) with
\(F(\gamma u)=\delta(\gamma)F(u)\).  Since \(|\delta(\gamma)|=1\),
\(\log|F|\) descends to a harmonic function on \(X\), hence is constant;
the open mapping theorem then makes \(F\) constant and forces
\(\delta=1\).  This check uses neither the higher-rank stability theorem
nor an assumption that all flat holomorphic bundles are stable.
\end{remark}
